\documentclass[11pt]{article}
\usepackage[margin=1in]{geometry}
\usepackage{amsmath,amssymb,amsthm}
\usepackage{graphicx}
\usepackage{booktabs}
\usepackage{hyperref}
\usepackage{natbib}
\usepackage{xcolor}
\usepackage{microtype}
\usepackage{algorithm}
\usepackage{algpseudocode}

\newcommand{\Gr}{\mathrm{Gr}}
\newcommand{\BN}{\mathrm{BN}}
\newcommand{\R}{\mathbb{R}}
\newcommand{\TODO}[1]{\textcolor{red}{[TODO: #1]}}
\newcommand{\CITE}[1]{\textcolor{red}{[CITE: #1]}}

\title{Causal Geometry Audit of RNA Foundation Model Representations:\\
  Bracket Norm, Grassmannian Transportability, and Weight Factorization\\
  on HTT mRNA for Huntington's Disease}

\author{Elliot Tower\\
  Mechanistic Research\\
  \texttt{elliot@elliottower.ai}}

\date{July 2026}

\begin{document}
\maketitle

\begin{abstract}
RNA-targeting therapeutics face two geometric confounds: distinguishing
genuine structural features from length-dependent artifacts, and
predicting whether drug mechanisms transfer across pathogenic mutations.
We apply three independent audits to RNA-FM, a 99M-parameter RNA
foundation model, using Huntingtin (HTT) mRNA as a case study.
%
First, bracket-norm correction reveals which embedding dimensions encode
intensive structural properties of the CAG repeat region versus
extensive properties that scale with sequence length.
%
Second, Grassmannian subspace analysis quantifies how the
representation subspace drifts as CAG repeat count increases from
wild-type (21 repeats) through the disease threshold (36+) to severe
expansions (80 repeats), yielding a transportability score for
antisense oligonucleotide design.
%
Third, singular value decomposition of RNA-FM's weight matrices shows
\TODO{effective rank results} across layers and projections, with
cross-projection subspace overlap suggesting a shared low-dimensional
factor structure amenable to mechanistic decomposition.
%
All analyses, code, and data are publicly available.
\end{abstract}

%% ────────────────────────────────────────────
\section{Introduction}
\label{sec:intro}

Huntington's disease (HD) is caused by expansion of a CAG trinucleotide
repeat in exon~1 of the Huntingtin gene (HTT), producing a
polyglutamine-expanded protein that aggregates and causes
neurodegeneration~\CITE{MacDonald et al.~1993, HTT Collaborative Group}.
Normal alleles carry 6--35 CAG repeats; 36--39 repeats confer incomplete
penetrance; 40+ repeats are fully penetrant~\CITE{Rubinsztein et al.~1996}.
No disease-modifying therapy exists. The primary therapeutic strategy
targets HTT mRNA directly, using antisense oligonucleotides (ASOs) to
reduce production of the toxic protein~\CITE{Tabrizi et al.~2019, tominersen trial}.

The central challenge for RNA-targeting drug design is \emph{transportability}:
a drug mechanism validated on one CAG repeat length must transfer to the
repeat lengths seen in patients. Clinical trials of tominersen
(an ASO targeting HTT mRNA) were halted partly because the drug reduced
both mutant and wild-type HTT with insufficient selectivity~\CITE{Tabrizi et al.~2022}.
Understanding which structural features of the CAG repeat region are
preserved across repeat lengths---and which degrade---would inform
the design of allele-selective therapies.

RNA foundation models offer a new lens on this problem. RNA-FM, a
99M-parameter BERT-style model pre-trained on 23 million non-coding RNA
sequences, produces 640-dimensional contextual embeddings for each
nucleotide position~\CITE{Chen et al.~2022, RNA-FM}. These embeddings
capture structural and functional properties learned from evolutionary
data, providing a representation space where geometric analyses can
detect genuine structural signals.

We apply three complementary audits from the causal geometry
framework~\CITE{Tower, bracket norm; Tower, Grassmannian cohort transfer}
to RNA-FM embeddings of the HTT CAG repeat region:

\begin{enumerate}
\item \textbf{Bracket-norm correction.} RNA embeddings are confounded
  by sequence length: longer sequences produce higher-norm
  representations due to positional encoding and attention aggregation
  effects. The bracket norm~\CITE{Tower, bracket norm} normalizes
  feature magnitudes by $\sqrt{n}$ to distinguish \emph{intensive}
  structural properties (genuine features of the RNA fold) from
  \emph{extensive} properties (artifacts scaling with input length).

\item \textbf{Grassmannian transportability.} For each CAG repeat count
  $r$, we compute the principal component subspace of RNA-FM embeddings
  on the repeat region, yielding a point on the Grassmannian
  $\Gr(k, 640)$. The geodesic distance between the wild-type subspace
  ($r = 21$) and a mutant subspace ($r > 36$) quantifies how much the
  representation structure degrades---a direct measure of drug
  mechanism transferability~\CITE{Tower, Grassmannian cohort transfer}.

\item \textbf{Weight factorization.} Singular value decomposition of
  RNA-FM's attention and MLP weight matrices reveals the intrinsic
  dimensionality of each projection. Cross-projection subspace overlap
  (shared singular directions between $W_Q$ and $W_K$, between $W_V$
  and $W_O$) indicates whether the model's internal structure admits a
  shared factor bank decomposition, as demonstrated for GPT-2 in prior
  work~\CITE{Tower \& Nanda, factorization circuits}.
\end{enumerate}

Together, these three audits provide a geometric validity check on any
structural claim derived from RNA-FM embeddings: features surviving
bracket-norm correction, Grassmannian transportability, and weight-space
factorization analysis carry higher evidential weight than features
identified by any single method.


%% ────────────────────────────────────────────
\section{Background}
\label{sec:background}

\subsection{RNA-FM architecture}

RNA-FM~\CITE{Chen et al.~2022} is a BERT-style~\CITE{Devlin et al.~2019}
transformer encoder with 12 layers, 20 attention heads per layer,
hidden dimension $d = 640$, and intermediate MLP dimension $d_{\mathrm{ff}} = 5120$.
The model uses a nucleotide-level vocabulary of 28 tokens (A, C, G, U,
plus special tokens) and was pre-trained with masked language modeling
on 23.7 million non-coding RNA sequences from RNAcentral~\CITE{RNAcentral}.
%
RNA-FM achieves state-of-the-art performance on RNA secondary structure
prediction, producing per-nucleotide embeddings that encode structural
context without explicit thermodynamic modeling.

\subsection{HTT mRNA and CAG repeat structure}

The human HTT mRNA (RefSeq NM\_002111.7) spans 13{,}669 nucleotides.
The pathogenic CAG repeat lies in exon~1, with normal alleles containing
approximately 21 repeats (63 nucleotides). The repeat region and its
flanking sequences create a structured RNA domain whose conformation
depends on repeat length: longer repeats form more stable hairpin
structures that may affect mRNA splicing, localization, and
translation~\CITE{de Mezer et al.~2011; Sobczak et al.~2003}.

\subsection{Bracket norm}

The bracket norm~\CITE{Tower, bracket norm} is a normalization that
distinguishes intensive from extensive properties in high-dimensional
feature spaces. For a feature vector $\mathbf{f} \in \R^d$ computed
from $n$ observations, the bracket-norm corrected value is
$\BN(\mathbf{f}) = \|\mathbf{f}\| / \sqrt{n}$.
%
Features whose bracket-norm corrected magnitude is invariant to $n$
represent genuine intensive properties; features whose corrected
magnitude varies with $n$ are extensive and may be confounded by
sampling density or input size. In the RNA context, $n$ is the
sequence length, and the correction identifies embedding dimensions
that encode per-nucleotide structural information rather than
global sequence statistics.

\subsection{Grassmannian subspace analysis}

A $k$-dimensional subspace of $\R^d$ is a point on the Grassmannian
manifold $\Gr(k, d)$~\CITE{Edelman et al.~1998}. Given two subspaces
$U, V \in \Gr(k, d)$, the principal angles
$\theta_1 \leq \cdots \leq \theta_k$ between them are computed via the
SVD of $U^\top V$, and the geodesic distance is
$d_g(U, V) = \sqrt{\sum_i \theta_i^2}$.
%
This distance quantifies structural similarity between representation
subspaces: small geodesic distance implies that the same geometric
features are active in both conditions, while large distance indicates
a qualitative shift in representation structure.
%
When applied to RNA embeddings across CAG repeat variants, this yields
a \emph{transportability score}: the degree to which drug-relevant
structural features identified on wild-type RNA persist in disease-length
expansions.


%% ────────────────────────────────────────────
\section{Methods}
\label{sec:methods}

\subsection{Data preparation}

We obtained the full-length HTT mRNA sequence (NM\_002111.7,
13{,}669~nt) from NCBI RefSeq. The native CAG repeat region was
identified by regex matching, yielding 21 tandem CAG repeats at
positions 148--211.

To study transportability across repeat lengths, we constructed
synthetic HTT mRNA variants by replacing the native CAG repeat with
$r \in \{15, 21, 27, 36, 40, 50, 60, 80\}$ tandem repeats, preserving
50~nt of flanking sequence on each side. This set spans
sub-normal (15), normal (21), intermediate (27), reduced-penetrance
threshold (36), full-penetrance (40), and severe expansion (50--80)
alleles.

\subsection{Embedding extraction}

We loaded RNA-FM weights from the HuggingFace
repository~\CITE{multimolecule/rnafm} into a BERT encoder using the
\texttt{transformers} library~\CITE{Wolf et al.~2020}, with
architecture parameters matched to the published
specification (12 layers, 640 hidden, 20 heads, 5120 intermediate).
%
Input sequences were tokenized at single-nucleotide resolution
(A$\to$5, C$\to$6, G$\to$7, U$\to$8) with CLS and EOS tokens.
Embeddings from the final hidden layer and all intermediate layers
were extracted for each variant.

\subsection{Bracket-norm analysis}

For windows of size $n \in \{100, 200, 300, 400, 500\}$~nt centered on
the CAG repeat region, we computed the mean embedding norm
$\bar{\ell}(n) = \frac{1}{n} \sum_{i=1}^{n} \|\mathbf{h}_i\|$
and the bracket-norm corrected value $\BN(n) = \bar{\ell}(n) / \sqrt{n}$.
%
An intensive property yields constant $\BN(n)$ across window sizes;
an extensive property yields $\BN(n)$ that decreases as $1/\sqrt{n}$.

We also computed per-dimension bracket norms: for each of the 640
embedding dimensions, we computed the $L^2$ norm across positions and
divided by $\sqrt{n}$, identifying the dimensions with the highest
intensive signal.

\subsection{Grassmannian transportability}

For each CAG repeat variant $r$, we computed RNA-FM embeddings on the
synthetic sequence and extracted the final-layer representations for
the repeat-plus-flanking region. PCA with $k = 10$ components yielded
a subspace $U_r \in \Gr(10, 640)$.

The geodesic distance $d_g(U_{21}, U_r)$ between the wild-type
subspace and each variant subspace quantifies transportability.
We also computed the full principal angle spectrum to identify which
subspace dimensions are preserved versus disrupted by repeat expansion.

\subsection{Layer-wise analysis}

For a 500~nt window centered on the CAG repeat, we extracted hidden
states from all 13 layers (embedding + 12 transformer layers) and
performed PCA on each. The effective dimensionality (number of
components for 95\% variance) characterizes how concentrated the
representation is at each processing stage.

We computed cross-layer subspace similarity as the mean cosine of
principal angles between 10-dimensional PCA subspaces of different
layers, yielding a $13 \times 13$ similarity matrix.

\subsection{Weight matrix factorization}

For each of the 12 transformer layers, we performed SVD on all six
projection matrices ($W_Q, W_K, W_V, W_O \in \R^{640 \times 640}$
and $W_{\mathrm{in}} \in \R^{5120 \times 640}$,
$W_{\mathrm{out}} \in \R^{640 \times 5120}$).
%
The effective rank at 95\% cumulative variance and the spectral decay
ratio $\sigma_1 / \sigma_{10}$ characterize the intrinsic
dimensionality of each projection.

To assess whether projections share a common factor structure, we
computed the subspace overlap between the top-20 right singular vectors
of each attention projection pair ($W_Q$--$W_K$, $W_Q$--$W_V$,
$W_K$--$W_V$, $W_Q$--$W_O$) within each layer. High overlap indicates
that the same input-space directions are important across projections,
supporting a shared factor bank decomposition.


%% ────────────────────────────────────────────
\section{Results}
\label{sec:results}

\TODO{Fill with actual numbers from analysis\_results\_latest.json once
the pipeline completes. The structure below matches the saved JSON fields.}

\subsection{Bracket-norm correction identifies intensive RNA features}

\TODO{Raw norms scale with window size. BN-corrected norms:
report whether they stabilize (intensive) or continue to decrease (extensive).
Report top intensive dimensions.}

% \begin{figure}[h]
% \centering
% \includegraphics[width=0.9\textwidth]{../figures/bracket_norm.png}
% \caption{(a) Mean embedding norm increases with window size, reflecting
% the extensive contribution of positional encoding. (b) Bracket-norm
% corrected values \TODO{stabilize/decrease}, indicating that
% \TODO{interpretation}.}
% \label{fig:bracket_norm}
% \end{figure}

\subsection{Grassmannian distance increases with CAG expansion}

\TODO{Report geodesic distances for each repeat count. Key question:
is the increase monotonic? Is there a phase transition near the
disease threshold (36 repeats)?}

% \begin{figure}[h]
% \centering
% \includegraphics[width=0.9\textwidth]{../figures/transportability.png}
% \caption{(a) Geodesic distance from the wild-type ($r=21$) subspace
% increases with CAG repeat count. \TODO{Note any nonlinearity near
% $r=36$.} (b) Principal angle heatmap shows which subspace dimensions
% are preserved (small angles) versus disrupted (large angles) by
% expansion.}
% \label{fig:transport}
% \end{figure}

\subsection{Layer-wise effective dimensionality}

\TODO{Report effective dims. Early layers typically high-dimensional,
later layers compress to task-relevant subspace.}

% \begin{figure}[h]
% \centering
% \includegraphics[width=0.9\textwidth]{../figures/cross_layer_similarity.png}
% \caption{Cross-layer subspace similarity matrix. \TODO{Describe
% block structure if present.}}
% \label{fig:layers}
% \end{figure}

\subsection{Weight matrices admit low-rank factorization}

\TODO{Report effective ranks. If attention matrices have effective
rank $\ll 640$, the shared factor bank decomposition is geometrically
motivated.}

% \begin{figure}[h]
% \centering
% \includegraphics[width=0.9\textwidth]{../figures/weight_factorization.png}
% \caption{(a) Effective rank (95\% variance) by layer and projection type.
% (b) Spectral decay ratio $\sigma_1/\sigma_{10}$ heatmap.}
% \label{fig:factorize}
% \end{figure}

\subsection{CAG repeat region is geometrically distinct}

\TODO{Report PCA separation, norm difference between CAG and flanking
positions.}

% \begin{figure}[h]
% \centering
% \includegraphics[width=0.9\textwidth]{../figures/position_analysis.png}
% \caption{(a) PCA of RNA-FM embeddings colored by region (CAG vs.\
% flanking). (b) Position-wise embedding norm shows \TODO{elevated/reduced}
% norm in the CAG region.}
% \label{fig:positions}
% \end{figure}


%% ────────────────────────────────────────────
\section{Discussion}
\label{sec:discussion}

\subsection{Implications for RNA-targeting drug design}

The transportability analysis provides a direct, quantitative answer to
a question facing ASO drug designers: how much does the structural
representation of HTT mRNA change between wild-type and disease-length
CAG expansions? \TODO{Interpret the geodesic distance curve.}

The bracket-norm analysis complements this by identifying which
embedding dimensions carry genuine structural information. Drug
design pipelines that screen RNA structural features should weight
intensive dimensions (high bracket-norm) over extensive ones to
avoid length-dependent artifacts.

\subsection{Weight factorization and mechanistic interpretability}

The SVD analysis reveals whether RNA-FM's internal computations are
amenable to the shared factor bank decomposition demonstrated for
GPT-2~\CITE{Tower \& Nanda, factorization circuits}. \TODO{Interpret
effective rank and cross-projection overlap results.}

If confirmed, a factorized RNA-FM would enable direct inspection of
which learned structural motifs (factors) are recruited by each
attention head and MLP projection, providing mechanistic interpretability
of RNA language model predictions. This addresses a key limitation
of current RNA structure prediction: models achieve high accuracy
without revealing which structural features drive their predictions.

\subsection{Limitations}

Several limitations qualify these results. RNA-FM was pre-trained
on non-coding RNA, not mRNA; its representations of protein-coding
sequences like HTT may not capture coding-specific structural features.
The PCA-based subspace extraction assumes linearity in the
representation geometry; nonlinear structure would require manifold
learning approaches. The synthetic repeat variants do not capture
somatic mosaicism or epigenetic modifications present in patient
tissue. The weight SVD characterizes the model's \emph{capacity}
for low-rank structure, but the full distillation-based factorization
(fitting a shared factor bank and per-projection selectors) requires
GPU training that we defer to subsequent work.

\subsection{Future directions}

Three extensions follow directly. First, cross-view validation using
thermodynamic structure prediction (EternaFold, RNAfold) as an
independent view, assessing whether the Grassmannian subspace identified
by RNA-FM aligns with physics-based structural ensembles. Second, full
distillation-based factorization of RNA-FM to extract interpretable
factors, following the pipeline established for
GPT-2~\CITE{Tower \& Nanda, factorization circuits}. Third, application
to additional RNA therapeutic targets (SARS-CoV-2 frameshifting element,
DMD exon-skipping targets, SMA survival motor neuron) to test whether
the bracket-norm and transportability metrics generalize across RNA
drug targets.


%% ────────────────────────────────────────────
\section{Conclusion}
\label{sec:conclusion}

We applied three geometric audits---bracket-norm correction,
Grassmannian transportability, and weight factorization---to RNA-FM
representations of the HTT CAG repeat region. \TODO{One-sentence
summary of key finding.} These tools provide a principled validity
check for RNA structural claims derived from foundation model
embeddings, and a quantitative transportability metric directly
relevant to RNA-targeting drug design for Huntington's disease.

All code, data, and analysis results are available at
\url{https://github.com/ElliottTower/causal-rna} under the MIT license.


%% ────────────────────────────────────────────
\bibliographystyle{plainnat}
% \bibliography{refs}

\begin{thebibliography}{99}

\bibitem[Chen et al.(2022)]{rnafm}
Chen, J., Hu, Z., Sun, S., Tan, Q., Wang, Y., Yu, Q., Zong, L., Hong, L., Xiao, J., King, I., et al.
\newblock Interpretable RNA foundation model from unannotated data for highly accurate RNA structure and function predictions.
\newblock \emph{arXiv:2204.00300}, 2022.

\bibitem[de~Mezer et al.(2011)]{demezer2011}
de~Mezer, M., Wojciechowska, M., Napierala, M., Sobczak, K., and Krzyzosiak, W.~J.
\newblock Mutant {CAG} repeats of {Huntingtin} transcript fold into hairpins, form nuclear foci and are targets for {RNA} interference.
\newblock \emph{Nucleic Acids Research}, 39(9):3852--3863, 2011.

\bibitem[Devlin et al.(2019)]{devlin2019bert}
Devlin, J., Chang, M.-W., Lee, K., and Toutanova, K.
\newblock {BERT}: Pre-training of deep bidirectional transformers for language understanding.
\newblock In \emph{NAACL-HLT}, 2019.

\bibitem[Edelman et al.(1998)]{edelman1998}
Edelman, A., Arias, T.~A., and Smith, S.~T.
\newblock The geometry of algorithms with orthogonality constraints.
\newblock \emph{SIAM Journal on Matrix Analysis and Applications}, 20(2):303--353, 1998.

\bibitem[MacDonald et al.(1993)]{macdonald1993}
MacDonald, M.~E., Ambrose, C.~M., Duyao, M.~P., et al.
\newblock A novel gene containing a trinucleotide repeat that is expanded and unstable on {Huntington's} disease chromosomes.
\newblock \emph{Cell}, 72(6):971--983, 1993.

\bibitem[Rubinsztein et al.(1996)]{rubinsztein1996}
Rubinsztein, D.~C., Leggo, J., Coles, R., et al.
\newblock Phenotypic characterization of individuals with 30--40 {CAG} repeats in the {Huntington} disease ({HD}) gene reveals {HD} cases with 36 repeats and apparently normal elderly individuals with 36--39 repeats.
\newblock \emph{American Journal of Human Genetics}, 59(1):16, 1996.

\bibitem[Sobczak et al.(2003)]{sobczak2003}
Sobczak, K., de~Mezer, M., Michlewski, G., Krol, J., and Krzyzosiak, W.~J.
\newblock {RNA} structure of trinucleotide repeats associated with human neurological diseases.
\newblock \emph{Nucleic Acids Research}, 31(19):5469--5482, 2003.

\bibitem[Tabrizi et al.(2019)]{tabrizi2019}
Tabrizi, S.~J., Leavitt, B.~R., Landwehrmeyer, G.~B., et al.
\newblock Targeting {Huntingtin} expression in patients with {Huntington's} disease.
\newblock \emph{New England Journal of Medicine}, 380(24):2307--2316, 2019.

\bibitem[Tabrizi et al.(2022)]{tabrizi2022}
Tabrizi, S.~J., Estevez-Fraga, C., van Roon-Mom, W.~M.~C., et al.
\newblock Potential disease-modifying therapies for {Huntington's} disease: lessons learned and future opportunities.
\newblock \emph{The Lancet Neurology}, 21(7):645--658, 2022.

\bibitem[Wolf et al.(2020)]{wolf2020transformers}
Wolf, T., Debut, L., Sanh, V., et al.
\newblock Transformers: State-of-the-art natural language processing.
\newblock In \emph{EMNLP: System Demonstrations}, 2020.

\end{thebibliography}

\end{document}
